\documentclass[10pt,a4paper]{amsart}
\usepackage[margin=19mm]{geometry}
\usepackage{amsmath,amssymb,mathtools}
\usepackage[scr=rsfs,cal=euler]{mathalfa}
\usepackage{xfrac}
\usepackage[unicode,hidelinks,bookmarks=false]{hyperref}
\newcommand{\ie}{i.e.\ }
\newcommand{\cf}{cf.\ }
\newcommand{\resp}{respectively\ }
\newcommand{\Subsection}{\subsection}
\providecommand{\nobreakdash}{\nobreak\hbox{-}}
\setlength{\emergencystretch}{2em}
\multlinegap0pt
\usepackage{amssymb}
\usepackage{enumerate}
\usepackage{xfrac}
\usepackage{caption}
\usepackage[normalem]{ulem}
\newtheorem{lemma}{Lemma}
\title{Explicit form of relaxation tensor for isotropic extended Burgers model and its spectral inversion}
\author{Youjun Deng, Ching-Lung Lin, Gen Nakamura}
\date{Source: arXiv:2602.10636v1 (CC BY 4.0)}
\begin{document}
\maketitle

\noindent\emph{Source and licence.} Explicit form of relaxation tensor for isotropic extended Burgers model and its spectral inversion. Version arXiv:2602.10636v1 (CC BY 4.0), distributed by arXiv under the Creative Commons Attribution 4.0 International licence, http://creativecommons.org/licenses/by/4.0/, as recorded in the arXiv licence metadata for this exact version and shown on the abstract page https://arxiv.org/abs/2602.10636v1. Full licence notice: \texttt{licenses/arxiv-2602.10636-CC-BY-4.0.txt}.\par\smallskip

\noindent\textit{Editorial scope.} Counted unit: the article's lemma s\_lem3.3 (source0.tex lines 1722--1725). The article's complete original proof of it is delivered with it (source0.tex lines 1726--1814, one \texttt{proof} environment of 89 lines). No mathematics is written, repaired or re-derived: the statement and the proof are the article's own lines, in the article's own notation, and no retained line is substituted or re-flowed.  No further source-local context is needed: every cross-reference of every retained line resolves inside the two delivered spans.  Nothing was omitted from any retained span, so the package carries no omission record.  No retained line contains a citation, so the wrapper carries no bibliography and no bibliography entry.  No retained line contains a figure environment, an image-inclusion command or drawing code, so no image is delivered and the package declares no asset dependency.  Editorial formatting only, in addition: an AMS \texttt{amsart} preamble that restates the article's own declarations and macros used by the retained lines, namely the environments \texttt{lemma} on separate counters, together with the standard packages those lines need.  The retained environments print in this document as Lemma 1 in order of appearance, whereas the article prints the same items with the numbering of its own full document.  The complete native source of the article as distributed in the arXiv e-print for this exact version is bundled unchanged as \texttt{sources/194414/source0.tex}.
\begin{lemma}\label{s_lem3.3}
There exists a unique $r_\ell\in((\ell-1)\pi, \ell \pi)$ such that $f(r_\ell)=0$ for each  $\ell=1,2,\cdots$. In fact, we have
$r_l\in((l-\frac{1}{2})\pi, l\pi)$ for each  $l=1,2,\cdots$. 
\end{lemma}

\begin{proof}
By direct computations, we have
\begin{equation}\label{s3.18}
\begin{aligned}
f'(\eta)=(\lambda_0+2\mu_0)\eta^2\cos\eta+2\lambda_0\eta\sin\eta
\end{aligned}
\end{equation}
and
\begin{equation}\label{s3.19}
\begin{aligned}
f''(\eta)=4(\lambda_0+\mu_0)\cos\eta+\big(2\lambda_0-(\lambda_0+2\mu_0)\eta^2\big)\sin\eta.
\end{aligned}
\end{equation}
We divide $\ell$ into two cases.

\medskip
\noindent
Case 1: $\ell=2m-1$ with a positive integer $m$.

From \eqref{s3.18}, we have that 
\begin{equation}\label{s3.20}
\begin{array}{l}
\begin{cases}
f'(\eta)>0,&\mbox{\rm if}\quad \eta\in\big((2m-2)\pi, (2m-\frac{3}{2})\pi\big),\\
f'(\eta)<0    & {\rm if}\quad \eta=(2m-1)\pi.
\end{cases}
\end{array}
\end{equation}
Since
\begin{equation*}
\begin{aligned}
2\lambda_0-(\lambda_0+2\mu_0)\eta^2<0\,\,\text{for}\,\,\eta\in[(2m-\frac{3}{2})\pi,(2m-1)\pi]
\end{aligned}
\end{equation*}
if $\eta\in[(2m-\frac{3}{2})\pi,(2m-1)\pi]$. Thus, we obtain that
\begin{equation}\label{s3.21}
\begin{aligned}
f''(\eta)<0
\end{aligned}
\end{equation}
if $\eta\in\big((2m-\frac{3}{2})\pi,(2m-1)\pi)\big)$. 
By \eqref{s3.20} and \eqref{s3.21}, there exists a unique $\eta_{2m-1}\in \big((2m-\frac{3}{2})\pi,(2m-1)\pi\big)$ such that $f'(\eta_{2m-1})=0$. We also have that
\begin{equation}\label{s3.22}
\begin{array}{l}
\begin{cases}
f'(\eta)>0,&\mbox{\rm if}\quad \eta\in((2m-2)\pi, \eta_{2m-1}),\\
f'(\eta)<0    & {\rm if}\quad \eta\in(\eta_{2m-1}, (2m-1)\pi).
\end{cases}
\end{array}
\end{equation}
Hence, by $f((2m-2)\pi)> 0$, $f((2m-1)\pi)< 0$ and \eqref{s3.22}, there exists a unique $r_{2m-1}\in (\eta_{2m-1},(2m-1)\pi)$ such that $f(r_{2m-1})=0$.

\medskip\noindent
Case 2: $\ell=2m$ with a positive integer $m$.

From \eqref{s3.18}, we have 
\begin{equation}\label{s3.23}
\begin{array}{l}
\begin{cases}
f'(\eta)<0,&\mbox{\rm if}\quad \eta\in\big((2m-1)\pi, (2m-\frac{1}{2})\pi\big),\\
f'(\eta)>0    & {\rm if}\quad \eta=2m\pi.
\end{cases}
\end{array}
\end{equation}
Since
\begin{equation*}
\begin{aligned}
2\lambda_0-(\lambda_0+2\mu_0)\eta^2<0\,\,\,\,\text{for}\,\,\,\,
\eta\in[(2m-\frac{1}{2})\pi,2m\pi],
\end{aligned}
\end{equation*}
we have
\begin{equation}\label{s3.24}
\begin{aligned}
f''(\eta)>0\,\,\,\,\text{for}\,\,\,\,\eta\in[(2m-\frac{1}{2})\pi,2m\pi)].
\end{aligned}
\end{equation}
By \eqref{s3.23} and \eqref{s3.24}, there exists a unique $\eta_{2m}\in \big((2m-\frac{1}{2})\pi,2m\pi\big)$ such that $f'(\eta_{2m})=0$. We also have
\begin{equation}\label{s3.25}
\begin{array}{l}
\begin{cases}
f'(\eta)<0,&\mbox{\rm if}\quad \eta\in((2m-1)\pi, \eta_{2m}),\\
f'(\eta)>0    & {\rm if}\quad \eta\in(\eta_{2m}, 2m\pi).
\end{cases}
\end{array}
\end{equation}
Hence, by $f((2m-1)\pi)< 0$, $f(2m\pi)> 0$ and \eqref{s3.25}, there exists a unique $r_{2m}\in (\eta_{2m},2m\pi)$ such that $f(r_{2m})=0$.

\end{proof}

\end{document}
